\begin{proof}[Proof of \cref{obj:lem:markov-polynomial}]
\begin{enumerate}
\item For a continuous real function \(f\) on \([-1,1]\), define
\[
\|f\|_{\infty}
  :=\sup\{|f(x)|:x\in[-1,1]\}.
\]
The set in this definition is nonempty and bounded, by compactness and continuity. Consequently, for every \(c\in\mathbb R\),
\[
\|f\|_{\infty}\le c
\quad\Longleftrightarrow\quad
|f(x)|\le c\quad\text{for every }x\in[-1,1].
\]
Indeed, each value is bounded by the supremum, and any common upper bound bounds the supremum. We will establish
\[
\|g'\|_{\infty}\le L^2\|g\|_{\infty}
\]
and then apply this equivalence to the hypothesis and conclusion.
% lean: Causalean.Mathlib.Analysis.FinitePolynomialAlternationDuality.intervalSupNorm_le_iff

\item If \(L=0\), the polynomial \(g\) is constant, so
\[
g'=0,\qquad \|g'\|_{\infty}=0=L^2\|g\|_{\infty}.
\]
Suppose henceforth that \(L>0\). Define the Chebyshev polynomials, viewed over either \(\mathbb R\) or \(\mathbb C\), by
\[
T_0(z)=1,\qquad T_1(z)=z,\qquad
T_{h+1}(z)=2zT_h(z)-T_{h-1}(z)\quad(h\ge1),
\]
and define their extrema and root angles by
\[
\xi_i:=\cos\frac{i\pi}{L}\quad(0\le i\le L),\qquad
\alpha_k:=\frac{(2k+1)\pi}{2L},\qquad
\rho_k:=\cos\alpha_k\quad(0\le k<L).
\]
The extrema are distinct and decreasing, and
\[
\xi_i\in[-1,1],\qquad
T_L(\xi_i)=(-1)^i.
\]
In particular,
\[
|g(\xi_i)|\le\|g\|_{\infty},
\qquad
0\le |g(\xi_0)|\le\|g\|_{\infty}.
\]
For \(0\le i\le L\), define the cardinal polynomials
\[
\lambda_i(z):=
\prod_{\substack{0\le j\le L\\j\ne i}}
\frac{z-\xi_j}{\xi_i-\xi_j}.
\]
Lagrange interpolation and differentiation give, for every real \(x\),
\[
g(z)=\sum_{i=0}^{L}g(\xi_i)\lambda_i(z),
\qquad
g'(x)=\sum_{i=0}^{L}g(\xi_i)\lambda_i'(x).
\]
Thus, for \(x\in[-1,1]\),
\[
|g'(x)|
\le\sum_{i=0}^{L}|g(\xi_i)|\,|\lambda_i'(x)|
\le\|g\|_{\infty}\sum_{i=0}^{L}|\lambda_i'(x)|.
\]
It remains to bound the final sum by \(L^2\).
% lean: Causalean.Mathlib.Analysis.FinitePolynomialAlternationDuality.duffinSchaeffer_derivative_le

\item Fix \(x\in[-1,1]\), and define
\[
\varepsilon_{x,i}:=
\begin{cases}
-1,&\lambda_i'(x)<0,\\
1,&\lambda_i'(x)\ge0,
\end{cases}
\qquad
g_{\mathrm{sgn},x}(z):=\sum_{i=0}^{L}\varepsilon_{x,i}\lambda_i(z).
\]
These definitions include the case \(\lambda_i'(x)=0\). Interpolation gives
\[
\deg g_{\mathrm{sgn},x}\le L,\qquad
g_{\mathrm{sgn},x}(\xi_i)=\varepsilon_{x,i},
\qquad
|g_{\mathrm{sgn},x}(\xi_i)|\le1.
\]
We first establish its derivative bound at the roots of \(T_L\).

Define
\[
\kappa_L:=L\,2^{L-1},\qquad
\mathcal N_L(z):=\prod_{j=0}^{L}(z-\xi_j),\qquad
\eta_i:=
\left(\prod_{\substack{0\le j\le L\\j\ne i}}
(\xi_i-\xi_j)\right)^{-1}.
\]
The following polynomial identities will be used:
\[
\kappa_L\mathcal N_L(z)=(z^2-1)T_L'(z),
\qquad
(1-z^2)T_L''(z)-zT_L'(z)+L^2T_L(z)=0.
\]
For the first identity, both sides have degree \(L+1\), leading coefficient \(\kappa_L\), and vanish at every \(\xi_i\): the endpoints annul \(z^2-1\), and the interior extrema annul \(T_L'\). The second is the Chebyshev differential identity, obtained by differentiating
\(T_L(\cos\theta)=\cos(L\theta)\).

Let \(\zeta\in\mathbb R\) satisfy \(T_L(\zeta)=0\). The explicit roots \(\rho_k\) give
\[
|\zeta|<1,\qquad \zeta\ne\xi_i\quad(0\le i\le L).
\]
Hence \(\mathcal N_L(\zeta)\ne0\), and the nodal identity implies
\(T_L'(\zeta)\ne0\). At this root, the differential identity yields
\[
(1-\zeta^2)T_L''(\zeta)=\zeta T_L'(\zeta).
\]
Differentiating
\[
\kappa_L\prod_{j\ne i}(z-\xi_j)
=\frac{(z^2-1)T_L'(z)}{z-\xi_i}
\]
at \(\zeta\), and using the preceding identity, gives
\[
\lambda_i'(\zeta)
=
\eta_i\,
\frac{T_L'(\zeta)(1-\zeta\xi_i)}
     {\kappa_L(\zeta-\xi_i)^2}.
\]
Since the extrema decrease,
\[
(-1)^i\eta_i>0.
\]
Also \(1-\zeta\xi_i>0\), because \(|\zeta|<1\) and \(|\xi_i|\le1\). Therefore
\[
(-1)^i\lambda_i'(\zeta)T_L'(\zeta)>0.
\]
Interpolation of \(T_L\) gives
\[
T_L'(\zeta)=\sum_{i=0}^{L}(-1)^i\lambda_i'(\zeta).
\]
If \(T_L'(\zeta)>0\), all summands on the right are positive; if \(T_L'(\zeta)<0\), all are negative. In either case,
\[
\sum_{i=0}^{L}|\lambda_i'(\zeta)|=|T_L'(\zeta)|.
\]
Consequently,
\[
\begin{aligned}
|g_{\mathrm{sgn},x}'(\zeta)|
&=\left|\sum_{i=0}^{L}
g_{\mathrm{sgn},x}(\xi_i)\lambda_i'(\zeta)\right|\\
&\le\sum_{i=0}^{L}
|g_{\mathrm{sgn},x}(\xi_i)|\,|\lambda_i'(\zeta)|\\
&\le\sum_{i=0}^{L}|\lambda_i'(\zeta)|
=|T_L'(\zeta)|.
\end{aligned}
\]
% lean: Causalean.Mathlib.Analysis.FinitePolynomialAlternationDuality.abs_eval_derivative_le_chebyshev_at_root

\item At the fixed point \(x\), differentiation of the signed interpolant gives
\[
g_{\mathrm{sgn},x}'(x)
=\sum_{i=0}^{L}\varepsilon_{x,i}\lambda_i'(x)
=\sum_{i=0}^{L}|\lambda_i'(x)|.
\]
Indeed, a negative weight is multiplied by \(-1\), and a nonnegative weight by \(1\). We next propagate the root bounds from the preceding step to this derivative value.
% lean: Causalean.Mathlib.Analysis.FinitePolynomialAlternationDuality.sum_abs_chebyshevDerivativeWeight_le_sq

\item Define the cardinal polynomials at the roots and a reflected-root polynomial by
\[
\lambda_{\mathrm{root},k}(z):=
\prod_{\substack{0\le j<L\\j\ne k}}
\frac{z-\rho_j}{\rho_k-\rho_j}
\quad(0\le k<L),
\qquad
\mathcal R_x(z):=
2^{L-1}\prod_{k=0}^{L-1}(z+|x-\rho_k|).
\]
Empty products equal \(1\), so these definitions also apply when \(L=1\).
Since
\[
T_L(z)=2^{L-1}\prod_{k=0}^{L-1}(z-\rho_k),
\]
we have
\[
T_L'(\rho_k)\lambda_{\mathrm{root},k}(x)
=
2^{L-1}\prod_{\substack{0\le j<L\\j\ne k}}(x-\rho_j).
\]
The derivative \(g_{\mathrm{sgn},x}'\) has degree at most \(L-1\). Interpolating it at the \(L\) roots and using the root bounds therefore gives
\[
\begin{aligned}
|g_{\mathrm{sgn},x}'(x)|
&\le\sum_{k=0}^{L-1}
|g_{\mathrm{sgn},x}'(\rho_k)|\,
|\lambda_{\mathrm{root},k}(x)|\\
&\le\sum_{k=0}^{L-1}
|T_L'(\rho_k)\lambda_{\mathrm{root},k}(x)|\\
&=2^{L-1}\sum_{k=0}^{L-1}
\prod_{\substack{0\le j<L\\j\ne k}}|x-\rho_j|
=|\mathcal R_x'(0)|.
\end{aligned}
\]
The last equality follows by differentiating the product defining
\(\mathcal R_x\); all its derivative summands at zero are nonnegative.
Moreover, for every \(y\in\mathbb R\),
\[
|\mathcal R_x(iy)|
=2^{L-1}\prod_{k=0}^{L-1}
\sqrt{(x-\rho_k)^2+y^2}
=|T_L(x+iy)|,
\qquad
\deg\mathcal R_x\le L.
\]
% lean: Causalean.Mathlib.Analysis.FinitePolynomialAlternationDuality.exists_chebyshevRootReflection

\item We now establish the vertical comparison
\[
|T_L(x+iy)|\le |T_L(1+iy)|
\qquad(x\in[-1,1],\ y\in\mathbb R).
\]
Define the squared chord function and the list of paired chord squares by
\[
d_{\mathrm{ch}}(\beta):=2-2\cos\beta
\quad(\beta\in\mathbb R),
\qquad
\mathscr D_L(\theta):=
\bigl(d_{\mathrm{ch}}(\theta+\alpha_k),
      d_{\mathrm{ch}}(\theta-\alpha_k)\bigr)_{k=0}^{L-1},
\]
where successive pairs are concatenated into a list of length \(2L\).
For an even list of nonnegative entries, define
\[
\mathcal P_t(v_1,\ldots,v_{2h})
:=\prod_{j=1}^{h}(v_{2j-1}v_{2j}+t)
\qquad(h\in\mathbb N,\ t\ge0),
\]
with value \(1\) for the empty list.

Sorting a list in decreasing order and pairing adjacent entries maximizes
\(\mathcal P_t\). To see this, for
\(v_1\ge v_2\ge v_3\ge v_4\ge0\), the two exchanges satisfy
\[
\begin{aligned}
&(v_1v_2+t)(v_3v_4+t)
 -(v_1v_3+t)(v_2v_4+t)
=t(v_1-v_4)(v_2-v_3)\ge0,\\
&(v_1v_2+t)(v_3v_4+t)
 -(v_1v_4+t)(v_2v_3+t)
=t(v_1-v_3)(v_2-v_4)\ge0.
\end{aligned}
\]
Thus the two largest entries can be paired together without decreasing the product. Removing that pair and inducting proves the assertion; the other factors are nonnegative.

For the fixed \(x\), define
\[
\theta_x:=\arccos x,\qquad
\delta_{\mathrm{ang}}:=\frac{\pi}{2L}.
\]
Choose \(k_{\mathrm{shift}}\in\mathbb Z\) nearest to \(L\theta_x/\pi\), and put
\[
\phi_x:=\theta_x-\frac{k_{\mathrm{shift}}\pi}{L},
\qquad
|\phi_x|\le\delta_{\mathrm{ang}}.
\]
The two lists \(\mathscr D_L(\theta_x)\) and
\(\mathscr D_L(\phi_x)\) have the same multiset of entries: their angles run through the odd multiples of \(\pi/(2L)\), shifted by the respective angle, and adding an integer multiple of \(\pi/L\) permutes these angles modulo \(2\pi\).

Define
\[
q_{k,+}:=d_{\mathrm{ch}}(\alpha_k+|\phi_x|),
\qquad
q_{k,-}:=d_{\mathrm{ch}}(\alpha_k-|\phi_x|)
\quad(0\le k<L).
\]
All the arguments lie in \([0,\pi]\). Since \(d_{\mathrm{ch}}\) increases on this interval and consecutive root angles differ by \(2\delta_{\mathrm{ang}}\),
\[
q_{k,-}\le q_{k,+},
\qquad
q_{k,+}\le q_{k+1,-}\quad(0\le k<L-1).
\]
Hence reversing the root blocks and placing \(q_{k,+}\) before \(q_{k,-}\) gives a decreasing list. Replacing \(\phi_x\) by \(|\phi_x|\) merely exchanges the two entries within a block, and reversing blocks preserves their paired product. The sorting inequality therefore implies
\[
\mathcal P_{4y^2}\bigl(\mathscr D_L(\theta_x)\bigr)
\le
\mathcal P_{4y^2}\bigl(\mathscr D_L(\phi_x)\bigr).
\]
The elementary chord identity
\[
d_{\mathrm{ch}}(\theta+\alpha_k)
d_{\mathrm{ch}}(\theta-\alpha_k)
=4(\cos\theta-\rho_k)^2
\]
turns this into
\[
4^L\left(\prod_{k=0}^{L-1}|x-\rho_k+iy|\right)^2
\le
4^L\left(\prod_{k=0}^{L-1}|\cos\phi_x-\rho_k+iy|\right)^2.
\]
The products are nonnegative, so
\[
\prod_{k=0}^{L-1}|x-\rho_k+iy|
\le
\prod_{k=0}^{L-1}|\cos\phi_x-\rho_k+iy|.
\]
Also
\[
\rho_k\le\cos\delta_{\mathrm{ang}}
\le\cos\phi_x\le1.
\]
Thus each factor on the right is bounded by its endpoint counterpart:
\[
|\cos\phi_x-\rho_k+iy|^2
=(\cos\phi_x-\rho_k)^2+y^2
\le(1-\rho_k)^2+y^2
=|1-\rho_k+iy|^2.
\]
Multiplying these bounds and then multiplying by \(2^{L-1}\) proves the required vertical comparison.
% lean: Causalean.Mathlib.Analysis.FinitePolynomialAlternationDuality.norm_eval_chebyshev_le_endpoint_vertical

\item Define
\[
\mathcal S(z):=T_L(1+z).
\]
Its degree and derivative at zero are
\[
\deg\mathcal S=L,\qquad \mathcal S'(0)=T_L'(1)=L^2>0,
\]
and its roots are \(\rho_k-1<0\). The preceding steps give
\[
|\mathcal R_x(iy)|=|T_L(x+iy)|
\le|T_L(1+iy)|=|\mathcal S(iy)|
\qquad(y\in\mathbb R).
\]
On the closed right half-plane, define
\[
\mathcal F_x(z):=\frac{\mathcal R_x(z)}{\mathcal S(z)}
\qquad(\operatorname{Re}z\ge0).
\]
The denominator has no zero there. This quotient is holomorphic in the open right half-plane, continuous on its closure, and, since
\(\deg\mathcal R_x\le\deg\mathcal S\),
\[
\mathcal F_x(z)=O(1)\quad\text{as }|z|\to\infty.
\]
In particular it has the growth bound required by the
Phragm\'en--Lindel\"of principle and is bounded along the positive real axis.
Its boundary values satisfy
\[
|\mathcal F_x(iy)|\le1.
\]
The right-half-plane Phragm\'en--Lindel\"of principle therefore gives
\[
|\mathcal F_x(z)|\le1
\qquad(\operatorname{Re}z\ge0).
\]

Suppose, for a contradiction, that
\(|\mathcal R_x'(0)|>|\mathcal S'(0)|\), and define
\[
\lambda_{\mathrm{der}}:=
\frac{\mathcal R_x'(0)}{\mathcal S'(0)},
\qquad
\mathcal E_x(z):=\mathcal R_x(z)-\lambda_{\mathrm{der}}\mathcal S(z).
\]
Then
\[
|\lambda_{\mathrm{der}}|>1,\qquad \mathcal E_x'(0)=0.
\]
For \(\operatorname{Re}z\ge0\), the equality \(\mathcal E_x(z)=0\) would imply
\(\mathcal F_x(z)=\lambda_{\mathrm{der}}\), contradicting the quotient bound.
Thus every root of \(\mathcal E_x\) lies in the open left half-plane.

If \(\mathcal E_x\) is nonconstant, the Gauss--Lucas theorem places every root of its derivative in the convex hull of its roots, hence in the open left half-plane. This contradicts \(\mathcal E_x'(0)=0\).
If \(\mathcal E_x\) is constant, define
\[
c_{\mathcal E}:=\mathcal E_x(0).
\]
Then
\[
\mathcal F_x(iy)
=\lambda_{\mathrm{der}}+\frac{c_{\mathcal E}}{\mathcal S(iy)}
\longrightarrow\lambda_{\mathrm{der}}
\quad\text{as }y\to+\infty,
\]
because \(\deg\mathcal S=L>0\). The boundary bound implies
\(|\lambda_{\mathrm{der}}|\le1\), again a contradiction. We conclude that
\[
|\mathcal R_x'(0)|\le|\mathcal S'(0)|=L^2.
\]
% lean: Causalean.Mathlib.Analysis.FinitePolynomialAlternationDuality.norm_eval_derivative_zero_le_chebyshev_endpoint_of_vertical

\item Combining the signed-interpolation identity, root reflection, and derivative comparison yields
\[
\sum_{i=0}^{L}|\lambda_i'(x)|
=g_{\mathrm{sgn},x}'(x)
\le |g_{\mathrm{sgn},x}'(x)|
\le|\mathcal R_x'(0)|
\le L^2.
\]
Returning to the interpolation bound for \(g\), we obtain
\[
|g'(x)|\le L^2\|g\|_{\infty}
\qquad(x\in[-1,1]),
\]
and hence
\[
\|g'\|_{\infty}\le L^2\|g\|_{\infty}.
\]
The hypothesis and the supremum characterization in the first step give
\[
\|g\|_{\infty}\le M_b.
\]
Since \(L^2\ge0\), multiplication and transitivity give
\[
\|g'\|_{\infty}
\le L^2\|g\|_{\infty}
\le L^2M_b.
\]
Finally, each derivative value is bounded by its supremum, so
\[
|g'(x)|\le L^2M_b
\qquad\text{for every }x\in[-1,1].
\]
% lean: CausalSmith.Stat.AnnotationRarearmFrontier.MarkovPolynomial
\end{enumerate}
\end{proof}
